\subsection{関係式を保ったまま微分する}
$x,y$ の関係が一つの式で与えられ，$y$ を $x$ の関数として明示していない場合がある。その式が実際に微分可能な関数 $y=y(x)$ を定めている範囲では，両辺を $x$ で微分できる。これを\textbf{陰関数の微分}という。$y$ も $x$ とともに変わるので，たとえば $(y^2)'=2yy'$，$(xy)'=y+xy'$ となる。

\begin{example}{二通りの表し方から接線を求める}
\begin{enumerate}[label=(\arabic*),leftmargin=*]
 \item $x=t^2+1$，$y=t^3-t$ で表される曲線の $t=1$ における接線を求めよ。
 \item 円 $x^2+y^2=25$ の点 $(3,4)$ における接線を求めよ。
\end{enumerate}
\end{example}
(1) では $t=1$ のとき $(x,y)=(2,0)$ である。また，
\[
 \frac{\dd x}{\dd t}=2t,\qquad
 \frac{\dd y}{\dd t}=3t^2-1,\qquad
 \left.\frac{\dd y}{\dd x}\right|_{t=1}
 =\frac{3-1}{2}=1.
\]
したがって接線は $y=x-2$ である。$t$ を消去しなくても，それぞれの変化率の比で傾きを求められた。

(2) では，点 $(3,4)$ の近くで上半円 $y=\sqrt{25-x^2}$ を考えれば，$y$ が微分可能な関数であることを確かめられる。その上で，元の関係式を微分すると
\[
 2x+2yy'=0,\qquad y'=-\frac xy\quad(y\ne0).
\]
点 $(3,4)$ での傾きは $-\frac34$ だから，接線は
\[
 y-4=-\frac34(x-3),\qquad 3x+4y=25
\]
である。円全体を一つの関数 $y=f(x)$ とみなしたわけではなく，指定点の近くの枝を使っている。$y=0$ の点では上の商は使えず，実際に接線は垂直になる。

これで，一点での変化率から，区間の増減や接線の近似までを調べる道具がそろった。次章では，細かい区間での寄与を足して定積分を作り，この微分の議論とつなげる。特に，平均値の定理から得た「導関数が同じなら差は定数」という結果が，積分の計算を支えることになる。

